% ============================================================
%  Chapter 3: Method
% ============================================================

\chapter{Method}
\label{ch:method}

This chapter describes the \modnet{} architecture and the evolutionary
algorithm that discovers its structure. We begin with an overview of
the design principles, then describe the node model, the modular
organization, the layer-free connectivity, the fitness function, and
the evolutionary operators for mutation, growth, and pruning.

% ------------------------------------------------------------
\section{Design Principles}
\label{sec:design-principles}
% ------------------------------------------------------------

\modnet{} is built on five design principles that distinguish it from
existing neuroevolutionary architectures.

\begin{enumerate}
  \item \textbf{No predefined layers.} There is no distinction between
    input, hidden, and output layers. Every node is a general-purpose
    computational unit that may receive input from any other node and
    may send output to any other node. Inputs and outputs are simply
    designated nodes; they are not structurally privileged.

  \item \textbf{Modular but not hierarchical.} Nodes are grouped into
    modules, but modules are not ordered, stacked, or nested. A module
    is simply a subset of nodes; the network has no global hierarchy.
    This is consistent with the principle of low coupling and high
    cohesion \citep{parnas1972modularity}.

  \item \textbf{Recurrence as a first-class primitive.} Any node may
    connect to itself or to any other node, regardless of index. The
    network is therefore a general directed graph, not a directed
    acyclic graph. Recurrence is not added as a special case; it is a
    natural consequence of unrestricted connectivity.

  \item \textbf{Growth through addition and pruning.} The network
    begins with a small number of nodes and grows as needed. New nodes
    are inserted into modules, and existing nodes are removed when they
    become inactive or isolated. Growth is bidirectional: the network
    may both increase and decrease in size.

  \item \textbf{Gradient-free learning.} No gradient information is
    used at any stage. All learning is performed by a compact
    evolutionary algorithm operating on the topology and weights
    simultaneously. This makes the architecture compatible with
    non-differentiable activation functions and avoids the vanishing
    gradient problem that affects deep recurrent networks.
\end{enumerate}

These principles are motivated by the observation that biological
neural networks exhibit none of the structural regularities that are
imposed on artificial networks by design. The brain has no layers, no
fixed topology, and no global error signal; it discovers its structure
through developmental and evolutionary processes
\citep{gershenson2010guiding, plantec2024evolving}. \modnet{} is an
attempt to capture these properties in a minimal computational model.

% ------------------------------------------------------------
\section{Node Model}
\label{sec:node-model}
% ------------------------------------------------------------

A node in \modnet{} is a threshold unit with a weighted sum of inputs.
Formally, let $N$ be the total number of nodes in the network. For a
node $i \in \{0, \ldots, N-1\}$, let $\mathcal{S}_i$ be the set of
source indices from which node $i$ receives input. Each source
$s \in \mathcal{S}_i$ may be either an external input (denoted by a
negative index $-k$ for input $k$) or another node (denoted by its
index $j \geq 0$). Node $i$ is also equipped with a bias $b_i$ and a
threshold parameter $\theta_i$.

At each time step $t$, the state of node $i$ is updated according to

\begin{equation}
  z_i^{(t)} = b_i + \sum_{s \in \mathcal{S}_i} w_{i,s} \, a_s^{(t)},
  \label{eq:node-pre-activation}
\end{equation}

\begin{equation}
  a_i^{(t+1)} =
  \begin{cases}
    1 & \text{if } z_i^{(t)} \geq 0, \\
    0 & \text{otherwise},
  \end{cases}
  \label{eq:node-activation}
\end{equation}

where $a_s^{(t)}$ is the activation of source $s$ at time $t$. External
inputs are held constant across time steps: for an input index $-k$,
$a_{-k}^{(t)} = x_k$ for all $t$.

The threshold at zero in Equation~\ref{eq:node-activation} is a
convenience; a non-zero threshold can be absorbed into the bias term.
The activation function is the Heaviside step function, which is
non-differentiable. This is why we use evolutionary rather than
gradient-based learning.

\begin{figure}[h]
  \centering
  \begin{tikzpicture}[
    node distance=1.2cm,
    every node/.style={font=\small},
    neuron/.style={circle, draw, minimum size=1cm, thick},
    input/.style={rectangle, draw, minimum size=0.7cm, fill=gray!20},
    arrow/.style={-{Latex[length=2mm]}, thick},
  ]
    % Inputs
    \node[input] (x1) at (-3, 1.5) {$x_1$};
    \node[input] (x2) at (-3, 0)   {$x_2$};
    \node[input] (x3) at (-3, -1.5) {$x_3$};

    % Node
    \node[neuron] (n) at (0, 0) {$i$};

    % Output
    \node[neuron] (out) at (2.5, 0) {$j$};

    % Edges
    \draw[arrow] (x1) -- node[above, sloped, font=\scriptsize] {$w_{i,1}$} (n);
    \draw[arrow] (x2) -- node[above, font=\scriptsize] {$w_{i,2}$} (n);
    \draw[arrow] (x3) -- node[below, sloped, font=\scriptsize] {$w_{i,3}$} (n);
    \draw[arrow] (n) -- (out);

    % Recurrent loop
    \draw[arrow] (n) to[out=135, in=45, looseness=8]
      node[above, font=\scriptsize] {$w_{i,i}$} (n);

    % Bias
    \node[font=\scriptsize, below=0.6cm of n] {$z_i = b_i + \sum_s w_{i,s} a_s$};
  \end{tikzpicture}
  \caption{Node model. Each node computes a weighted sum of its
    sources plus a bias, then applies a Heaviside step function. The
    self-loop (dashed arrow) indicates that the node may receive input
    from itself, which introduces recurrence.}
  \label{fig:node-model}
\end{figure}

% ------------------------------------------------------------
\section{Modular Structure}
\label{sec:modular-structure}
% ------------------------------------------------------------

Nodes are partitioned into $M$ modules, where $M$ is a fixed
hyperparameter at initialization but the \emph{sizes} of the modules
are not fixed. Let $m_i \in \{0, \ldots, M-1\}$ denote the module
assignment of node $i$, and let $n_m$ denote the number of nodes in
module $m$, so that $\sum_{m=0}^{M-1} n_m = N$.

The module assignment is used only for two purposes: (i) to control
where new nodes are inserted during growth, and (ii) to measure the
degree of cross-module integration in the evolved network. It does
\emph{not} restrict connectivity: a node in module $m$ may receive
input from any node in any module, and may send output to any node in
any module.

\subsection{Module Sizes}

At initialization, all modules have the same size: $n_m = n_0$ for all
$m$. During evolution, the sizes change through growth and pruning.
A new node is inserted into a module chosen by weighted sampling, where
the weight of module $m$ is inversely proportional to its current size:

\begin{equation}
  P(\text{insert into module } m) \propto \frac{1}{n_m + 1}.
  \label{eq:module-growth-priority}
\end{equation}

This has the effect of keeping module sizes roughly balanced, while
still allowing one module to grow faster than others if the task
demands it.

\subsection{Output Module}

One module is designated as the \emph{output module}. The network's
output is read from a subset of nodes within this module; specifically,
the first $n_{\text{out}}$ nodes of the output module (in index order)
provide the output bits. During pruning, the output module is protected
from shrinking below $n_{\text{out}}$ nodes.

\subsection{Measuring Modularity}

We measure the degree of modularity in an evolved network using the
ratio of cross-module connections to total active connections:

\begin{equation}
  Q = \frac{\left| \{(i,j) : m_i \neq m_j, |w_{i,j}| > \epsilon\} \right|}
           {\left| \{(i,j) : |w_{i,j}| > \epsilon\} \right|},
  \label{eq:modularity-metric}
\end{equation}

where $\epsilon$ is a small threshold (we use $\epsilon = 0.05$) below
which a weight is considered inactive. A value of $Q = 0$ indicates a
completely modular network (no cross-module connections), while $Q = 1$
indicates a fully integrated network. We report $Q$ for each evolved
network in Chapter~\ref{ch:results}.

\begin{figure}[h]
  \centering
  \begin{tikzpicture}[
    every node/.style={font=\small},
    node0/.style={circle, draw, fill=blue!30, minimum size=0.6cm},
    node1/.style={circle, draw, fill=red!30, minimum size=0.6cm},
    node2/.style={circle, draw, fill=green!30, minimum size=0.6cm},
    node3/.style={circle, draw, fill=orange!30, minimum size=0.6cm},
    arrow/.style={-{Latex[length=1.5mm]}, thick, gray!70},
    cross/.style={-{Latex[length=1.5mm]}, thick, black},
  ]
    % Module 0
    \node[node0] (a1) at (0, 2) {};
    \node[node0] (a2) at (0, 1) {};
    \node[node0] (a3) at (0, 0) {};

    % Module 1
    \node[node1] (b1) at (2.5, 2) {};
    \node[node1] (b2) at (2.5, 1) {};
    \node[node1] (b3) at (2.5, 0) {};

    % Module 2
    \node[node2] (c1) at (5, 2) {};
    \node[node2] (c2) at (5, 1) {};
    \node[node2] (c3) at (5, 0) {};

    % Module 3 (output)
    \node[node3] (d1) at (7.5, 1.5) {};
    \node[node3] (d2) at (7.5, 0.5) {};

    % Intra-module edges (light)
    \draw[arrow] (a1) -- (a2);
    \draw[arrow] (a2) -- (a3);
    \draw[arrow] (b1) -- (b2);
    \draw[arrow] (b2) -- (b3);
    \draw[arrow] (c1) -- (c2);
    \draw[arrow] (c2) -- (c3);

    % Cross-module edges (bold)
    \draw[cross] (a1) -- (b2);
    \draw[cross] (a3) -- (c1);
    \draw[cross] (b1) -- (c3);
    \draw[cross] (b3) -- (d1);
    \draw[cross] (c1) -- (d2);
    \draw[cross] (a2) to[out=-60, in=-120] (d1);
    \draw[cross] (c3) -- (d1);

    % Recurrent loop
    \draw[cross, dashed] (b2) to[out=135, in=45, looseness=8] (b2);

    % Labels
    \node[above=0.3cm of a1, font=\small] {Module 0};
    \node[above=0.3cm of b1, font=\small] {Module 1};
    \node[above=0.3cm of c1, font=\small] {Module 2};
    \node[above=0.3cm of d1, font=\small] {Module 3 (out)};

    % Legend
    \node[anchor=west, font=\scriptsize] at (0, -1.2) {%
      \tikz\draw[thick, gray!70] (0,0) -- (0.4,0);
      \; intra-module \quad
      \tikz\draw[thick, black] (0,0) -- (0.4,0);
      \; cross-module \quad
      \tikz\draw[thick, dashed, black] (0,0) -- (0.4,0);
      \; recurrent
    };
  \end{tikzpicture}
  \caption{Modular layer-free architecture. Nodes (circles) are
    grouped into modules (colors), but cross-module connections
    (bold arrows) are unrestricted. The dashed self-loop on a node in
    Module 1 illustrates recurrence. Cross-module edges may connect
    any pair of modules, including from the output module back to
    earlier modules.}
  \label{fig:modular-architecture}
\end{figure}

% ------------------------------------------------------------
\section{Layer-Free Connectivity}
\label{sec:layer-free}
% ------------------------------------------------------------

The central architectural choice in \modnet{} is the absence of layers.
A connection from node $i$ to node $j$ is permitted if and only if
$i \neq j$ is false---that is, self-connections are explicitly allowed.
Formally, the connectivity matrix $\mathbf{W} \in \mathbb{R}^{N \times N}$
has no structural zeros: any entry $W_{i,j}$ may be non-zero,
regardless of the module assignments $m_i$ and $m_j$.

However, connections from an external input are handled separately.
We use a separate input weight matrix $\mathbf{W}_{\text{in}} \in
\mathbb{R}^{N \times K}$, where $K$ is the number of external inputs.
A node $i$ may receive input from any subset of the $K$ inputs, and
this subset is determined by evolution.

\begin{remark}
The absence of layers has two important consequences. First, the
concept of ``depth'' does not apply: a signal may traverse the network
in cycles of arbitrary length, and information may flow in any
direction. Second, there is no natural topological ordering of nodes;
the index $i$ is purely a labeling convention and does not imply any
temporal or structural priority.
\end{remark}

To ensure that the network has well-defined behavior, we evaluate it
for a fixed number of time steps $T$, which is a hyperparameter. The
state of all nodes is initialized to zero, and after $T$ steps the
output is read from the output module. This is analogous to the
``unrolling'' of a recurrent network, but with the crucial difference
that the unrolled graph is not required to be acyclic.

% ------------------------------------------------------------
\section{Recurrence}
\label{sec:recurrence}
% ------------------------------------------------------------

Recurrence in \modnet{} arises from two sources: self-connections and
cycles of length greater than one. Both are permitted, and both may
coexist in the same network.

A self-connection on node $i$ is simply a non-zero entry $W_{i,i}$.
Such a connection allows the node to retain information across time
steps: if the node is active at time $t$ and its self-weight is
positive, it may remain active at time $t+1$ even if no other input
changes.

A cycle of length $\ell > 1$ is a sequence of nodes
$i_1 \to i_2 \to \cdots \to i_\ell \to i_1$ where each consecutive
pair has a non-zero connection. Cycles allow information to be
propagated, transformed, and returned to earlier parts of the network.

We quantify recurrence using the number of self-connections:

\begin{equation}
  R = \left| \{i : |W_{i,i}| > \epsilon\} \right|.
  \label{eq:recurrence-metric}
\end{equation}

We report $R$ for each evolved network in Chapter~\ref{ch:results}.
As we will show, $R$ is not imposed by the design; it emerges
spontaneously and correlates with the temporal complexity of the task.

% ------------------------------------------------------------
\section{Fitness Evaluation}
\label{sec:fitness}
% ------------------------------------------------------------

Let $\mathcal{D} = \{(\mathbf{x}^{(k)}, \mathbf{y}^{(k)})\}_{k=1}^{K}$
be a set of input--output pairs. For a given network, let
$\hat{\mathbf{y}}^{(k)}$ be the output produced when the network is
initialized to zero and run for $T$ time steps with input
$\mathbf{x}^{(k)}$.

For Boolean tasks, the fitness is the negative mean squared error
between the predicted scalar and the target scalar:

\begin{equation}
  F_{\text{bool}}(\modnet{}) = -\frac{1}{K} \sum_{k=1}^{K}
    \left( \hat{y}^{(k)} - y^{(k)} \right)^2,
  \label{eq:fitness-boolean}
\end{equation}

where $\hat{y}^{(k)}$ is obtained by decoding the output bits as an
integer and dividing by the maximum representable value.

For regression tasks, the fitness is the negative mean squared error
between the decoded output and the target value:

\begin{equation}
  F_{\text{reg}}(\modnet{}) = -\frac{1}{K} \sum_{k=1}^{K}
    \left( \text{decode}(\hat{\mathbf{y}}^{(k)}) - y^{(k)} \right)^2,
  \label{eq:fitness-regression}
\end{equation}

where $\text{decode}(\cdot)$ is the same decoding function used in the
Boolean case. The use of a single fitness function for both Boolean
and continuous tasks is deliberate: it emphasizes that the architecture
does not need to be adapted to the task type.

\begin{remark}
In both cases, the fitness is negative, so that a higher fitness is
always better, and the maximum attainable fitness is zero. This
convention simplifies the selection mechanism and avoids the need for
a separate termination criterion.
\end{remark}

% ------------------------------------------------------------
\section{Evolutionary Algorithm}
\label{sec:evolution}
% ------------------------------------------------------------

The evolutionary algorithm in \modnet{} is a compact generational
algorithm with four operators: mutation, crossover, growth, and
pruning. The algorithm maintains a population of $P$ networks, where
$P$ is a hyperparameter. At each generation, the population is
evaluated, sorted by fitness, and used to produce the next generation.

\begin{algorithm}[h]
  \caption{ModNet Evolution}
  \label{alg:evolution}
  \begin{algorithmic}[1]
    \State Initialize population $\mathcal{P} \gets \{N_1, \ldots, N_P\}$
      with random topologies.
    \State Evaluate $F(N_p)$ for all $p \in \{1, \ldots, P\}$.
    \For{$g = 1, 2, \ldots, G$}
      \State Sort $\mathcal{P}$ by fitness in descending order.
      \State $\mathcal{E} \gets$ top $\lfloor P/5 \rfloor$ networks.
      \State $\mathcal{P}' \gets \mathcal{E}$.
      \While{$|\mathcal{P}'| < P$}
        \State $r \gets \text{Uniform}(0, 1)$.
        \If{$r < p_{\text{cross}}$ and $|\mathcal{E}| \geq 2$}
          \State $N_{\text{new}} \gets \text{Crossover}(N_a, N_b)$
            for random $N_a, N_b \in \mathcal{E}$.
        \ElsIf{$r < 1 - p_{\text{fresh}}$}
          \State $N_{\text{new}} \gets \text{Mutate}(N_a)$
            for random $N_a \in \mathcal{E}$.
        \Else
          \State $N_{\text{new}} \gets$ random new network.
        \EndIf
        \State $\mathcal{P}' \gets \mathcal{P}' \cup \{N_{\text{new}}\}$.
      \EndWhile
      \If{$g \bmod g_{\text{birth}} = 0$}
        \State Apply \textsc{Birth} to top 3 networks.
      \EndIf
      \If{$g \bmod g_{\text{prune}} = 0$}
        \State Apply \textsc{Prune} to top $\lfloor P/5 \rfloor$ networks.
      \EndIf
      \State $\mathcal{P} \gets \mathcal{P}'$.
      \State Evaluate $F(N_p)$ for all $N_p \in \mathcal{P}$
        with fitness changed.
    \EndFor
    \State \Return best network in $\mathcal{P}$.
  \end{algorithmic}
\end{algorithm}

\subsection{Mutation}

The mutation operator modifies the weights, biases, and connectivity
of a network. For each mutation step, a node $i$ is selected uniformly
at random, and one of four operations is applied:

\begin{enumerate}
  \item With probability $0.40$, perturb the weight of a randomly
    selected incoming connection: $W_{i,j} \gets W_{i,j} +
    \mathcal{N}(0, \sigma)$.
  \item With probability $0.20$, perturb the bias:
    $b_i \gets b_i + \mathcal{N}(0, \sigma)$.
  \item With probability $0.20$, perturb the weight of an input
    connection: $W^{\text{in}}_{i,k} \gets W^{\text{in}}_{i,k} +
    \mathcal{N}(0, \sigma)$.
  \item With probability $0.20$, either remove an existing connection
    (if the node has more than one active incoming connection) or add
    a new connection with a weight drawn from $\mathcal{N}(0, 1)$.
\end{enumerate}

The mutation rate $\sigma$ is not fixed; it grows when the population
reaches a plateau, according to the schedule

\begin{equation}
  \sigma \gets \min(\sigma_{\max}, \sigma \cdot \gamma)
  \quad \text{when } g_{\text{plateau}} > g_{\text{thresh}},
  \label{eq:sigma-adaptation}
\end{equation}

where $g_{\text{plateau}}$ counts generations since the last fitness
improvement, and $g_{\text{thresh}}$ is a hyperparameter. This allows
the population to escape local optima by increasing the exploration
radius when needed.

\subsection{Crossover}

The crossover operator combines two parent networks into a single
child. It requires that the two parents have the same number of nodes
and the same module sizes; if this is not the case, the operator
returns \texttt{None} and the caller falls back to mutation. For each
node $i$, with probability $0.5$, the entire row of $\mathbf{W}$,
the entire row of $\mathbf{W}^{\text{in}}$, and the bias $b_i$ are
copied from the second parent; otherwise they are kept from the first
parent.

This row-wise crossover is a lightweight alternative to the historical
markings used in NEAT \citep{stanley2002evolving}, which would require
tracking the lineage of each connection. Because \modnet{} does not
use historical markings, we restrict crossover to parents of identical
shape. This is a deliberate simplification: it reduces bookkeeping and
ensures that the child is always a valid network.

\subsection{Growth (Birth)}

The birth operator adds a single new node to the network. The target
module is chosen by weighted sampling according to
Equation~\ref{eq:module-growth-priority}, and the new node is inserted
at a position determined by its module assignment. The new node
receives a small number of random incoming connections, with weights
drawn from $\mathcal{N}(0, 1)$, and a random bias.

Growth is applied periodically (every $g_{\text{birth}}$ generations)
to the top three networks in the population. This ensures that the
network size is allowed to increase when the task demands it, but
only in response to sustained selection pressure.

\subsection{Pruning}

The pruning operator removes nodes that are no longer contributing to
the network's behavior. A node is considered \emph{dead} if its
activity across all training samples is below a lower threshold or
above an upper threshold:

\begin{equation}
  \bar{a}_i < \alpha_{\text{lo}} \quad \text{or} \quad
  \bar{a}_i > \alpha_{\text{hi}},
  \label{eq:dead-node}
\end{equation}

where $\bar{a}_i$ is the mean activation of node $i$ across all
samples and all time steps, and $\alpha_{\text{lo}}, \alpha_{\text{hi}}$
are hyperparameters. A node is also considered dead if it has neither
incoming nor outgoing active connections (an isolated node).

Pruning is applied every $g_{\text{prune}}$ generations to the top
$\lfloor P/5 \rfloor$ networks. To avoid destroying functional
networks, pruning is only accepted if the resulting network has
fitness at least as high as the original (up to a small tolerance
$\epsilon_{\text{prune}}$). This is the same ``safe prune'' strategy
used in the final version of our algorithm, and it ensures that
pruning never degrades performance.

\begin{figure}[h]
  \centering
  \begin{tikzpicture}[
    every node/.style={font=\small},
    alive/.style={circle, draw, fill=blue!30, minimum size=0.6cm},
    dead/.style={circle, draw, fill=gray!30, minimum size=0.6cm,
      dashed},
    arrow/.style={-{Latex[length=1.5mm]}, thick},
  ]
    % Before
    \node[alive] (a1) at (0, 2) {};
    \node[dead]  (a2) at (0, 1) {};
    \node[alive] (a3) at (0, 0) {};
    \node[alive] (b1) at (2, 2) {};
    \node[alive] (b2) at (2, 1) {};
    \node[dead]  (b3) at (2, 0) {};

    \draw[arrow] (a1) -- (a3);
    \draw[arrow] (a1) -- (b2);
    \draw[arrow] (a3) -- (b2);
    \draw[arrow] (b1) -- (a1);
    \draw[arrow] (b2) -- (a3);

    \node[font=\small] at (1, -1) {Before pruning};

    % Arrow
    \draw[-{Latex[length=3mm]}, thick] (3.5, 1) -- (5, 1);

    % After
    \node[alive] (c1) at (6.5, 2) {};
    \node[alive] (c3) at (6.5, 0.5) {};
    \node[alive] (d1) at (8.5, 2) {};
    \node[alive] (d2) at (8.5, 1) {};

    \draw[arrow] (c1) -- (c3);
    \draw[arrow] (c1) -- (d2);
    \draw[arrow] (c3) -- (d2);
    \draw[arrow] (d1) -- (c1);
    \draw[arrow] (d2) -- (c3);

    \node[font=\small] at (7.5, -1) {After pruning};
  \end{tikzpicture}
  \caption{Pruning removes nodes whose activity is either below
    $\alpha_{\text{lo}}$ or above $\alpha_{\text{hi}}$ across all
    training samples (dashed gray nodes). The remaining nodes and
    connections are re-indexed to form a compact network.}
  \label{fig:pruning}
\end{figure}

\subsection{Selection}

Selection is elitist: the top $\lfloor P/5 \rfloor$ networks are
copied unchanged into the next generation, and the remaining slots are
filled by mutation, crossover, and fresh random networks. The
proportion of fresh random networks is controlled by
$p_{\text{fresh}}$, which we set to $0.05$. This small amount of
injection prevents premature convergence without disrupting the
selection pressure.

% ------------------------------------------------------------
\section{Complexity Analysis}
\label{sec:complexity}
% ------------------------------------------------------------

The computational cost of one evaluation of a network with $N$ nodes
and $T$ time steps is $O(N^2 T)$ in the worst case, because each node
may receive input from any other node. In practice, the connectivity
is sparse (Section~\ref{sec:sparse}), so the cost is $O(E T)$, where
$E$ is the number of active connections. For the networks in our
experiments, $N \leq 30$ and $E \leq 300$, so a single evaluation
takes on the order of microseconds on a modern CPU.

The evolutionary algorithm performs $O(P G)$ evaluations, where $P$
is the population size and $G$ is the number of generations. With
$P = 200$ and $G = 1500$, this is $3 \times 10^5$ evaluations, which
takes approximately one minute on a single phone CPU.

% ------------------------------------------------------------
\section{Summary of Hyperparameters}
\label{sec:hyperparameters}
% ------------------------------------------------------------

Table~\ref{tab:hyperparameters} lists all hyperparameters used in our
experiments. We report the value used for most tasks; for tasks with
different values, the specific choice is noted in
Chapter~\ref{ch:experiments}.

\begin{table}[h]
  \centering
  \small
  \caption{Hyperparameters of \modnet{}.}
  \label{tab:hyperparameters}
  \begin{tabular}{lll}
    \toprule
    \textbf{Symbol} & \textbf{Description} & \textbf{Value} \\
    \midrule
    $M$ & Number of modules & 4 \\
    $n_0$ & Initial nodes per module & 6 \\
    $T$ & Time steps per evaluation & 4--8 \\
    $P$ & Population size & 200--300 \\
    $G$ & Generations & 400--1500 \\
    $p_{\text{cross}}$ & Crossover rate & 0.30 \\
    $p_{\text{fresh}}$ & Fresh random rate & 0.05 \\
    $\sigma_{\text{base}}$ & Base mutation scale & 0.30 \\
    $\sigma_{\max}$ & Maximum mutation scale & 1.50 \\
    $\gamma$ & Mutation growth factor & 1.25 \\
    $g_{\text{thresh}}$ & Plateau threshold & 25 \\
    $g_{\text{birth}}$ & Birth interval & 80--100 \\
    $g_{\text{prune}}$ & Prune interval & 15--20 \\
    $\epsilon_{\text{prune}}$ & Prune tolerance & $-0.001$ \\
    $\alpha_{\text{lo}}$ & Lower activity threshold & 0.02 \\
    $\alpha_{\text{hi}}$ & Upper activity threshold & 0.98 \\
    $n_{\text{out}}$ & Number of output nodes & 1 or 8 \\
    \bottomrule
  \end{tabular}
\end{table}

% ------------------------------------------------------------
\section{Reproducibility}
\label{sec:reproducibility}
% ------------------------------------------------------------

The complete implementation of \modnet{}, including the evolutionary
algorithm, the fitness functions, the mutation and pruning operators,
and the experimental framework that generates all figures and tables
reported in this thesis, is available as open-source software
\citep{modnet2026code}. The implementation depends only on NumPy and
the Python standard library; it requires no GPU and runs on commodity
hardware, including mobile phones.

All experiments are seeded, and the seed is reported in the
experimental setup (Chapter~\ref{ch:experiments}). The output of each
experiment is stored as JSON, CSV, SVG, and plain-text files, which
allows full reproduction and inspection of every evolved network.